\documentclass[12pt]{article}

%\usepackage[cm]{fullpage}
\setlength{\topmargin}{-1in}
\setlength{\oddsidemargin}{-0.25in}
\setlength{\evensidemargin}{-0.25in}
\setlength{\textwidth}{7in}
\setlength{\textheight}{10in}

\usepackage{array}
\usepackage{enumitem}
\usepackage{amssymb}
\usepackage{amsmath}
\usepackage{epstopdf}
\usepackage[T1]{fontenc}
\usepackage{graphics}
\usepackage{graphicx}
\usepackage{hyperref}
\usepackage{latexsym}
\usepackage{mathrsfs}
\usepackage{multicol}
\usepackage{setspace}
\usepackage{subfigure}
\usepackage{time}
\usepackage{times}

\newcommand{\To}{\Rightarrow}
\newcommand{\del}{\nabla}
\newcommand{\R}{\mathbb{R}}
\newcommand{\C}{\mathcal{C}}
\newcommand{\N}{\mathcal{N}}
\newcommand{\A}{\mathcal{A}}
\renewcommand{\O}{\mathcal{O}}
% \renewcommand{\N}{\mathbb{N}}
\newcommand{\ben}{\begin{enumerate}}
\newcommand{\een}{\end{enumerate}}
\newcommand{\eps}{\varepsilon}
\renewcommand{\epsilon}{\varepsilon}
\newcommand{\p}{\partial}
\newcommand{\<}{\langle}
\renewcommand{\>}{\rangle}
\renewcommand{\subset}{\subseteq}
\newcommand{\cont}{\Rightarrow \Leftarrow}
\newcommand{\back}{\backslash}
\newcommand{\norm}[1]{\|{#1}\|}
\newcommand{\abs}[1]{|{#1}|}
\newcommand{\ip}[1]{\langle{#1}\rangle}
\newcommand{\m}[1]{{\mathbf{#1}}}
\newcommand{\sts}{\m{S}^\textrm{T}\m{S}}
\newcommand{\rank}{\textrm{rank}}
\newcommand{\X}{\mathcal{X}}
\newcommand{\Y}{\mathcal{Y}}
\def\ds{\displaystyle}
\def\d{\partial}
\renewcommand{\arraystretch}{1.5}
\renewcommand{\L}{\mathscr{L}}

%\includegraphics[scale = 0.65, trim = 1in 3.25in 1in 3in]{blank_axes.pdf}


\pagestyle{empty}
\begin{document}
\noindent\rule{7in}{0.4pt}\\
\begin{center}
{\sc Differential Equations\\
Solutions of DEs}
\end{center}
\noindent\rule{7in}{0.4pt}
%%%%%%%%%%%%%%%%%%%%%%%%%%%%%%%%%%%%%%%%%%%%%%%%
\section*{Solution to a DE}
\doublespacing

\noindent  It might be helpful to make an analogy here by thinking about what it means to be a solution
to an algebraic equation.  Consider the following algebraic equation.
	\begin{equation} x^5 + 12x^2 +10x = 7x^3 + 24  \label{alg_eqn} \end{equation}
\noindent  Think about this for a moment:  what does it mean to be a solution to this equation?  We are often
asked to solve for things and we get an answer and move one, but what does it really mean in this case?
Could you verify whether something is or is not a solution?\\

\noindent  In the context of this algebraic equation, a solution would be an $x$-value that ``works" when we substitute
it back into the equation.  

\noindent  {\bf Algebraic Example A:}  Is $x = -1$ a solution to equation (\ref{alg_eqn}) above?
\newpage

%If we substitute $x = -1$ into each side of the equation, we have
%	\begin{align*}
%	\mbox{left hand side of the equation (LHS)}	%&= x^5 + 12x^2 +10x,\\
%								&= (-1)^5 + 12(-1)^2 +10(-1),\\
%								&= -1 + 12 - 11,\\
%								&= 0,\\
%	\end{align*}
%and
%	\begin{align*}
%	\mbox{right hand side of the equation (RHS)}	%&= 7x^3 + 24,\\
%									&= 7(-1)^3 + 24,\\
%									&= -7 + 24,\\
%									&= 17.
%	\end{align*}
%\noindent  When we substitute $x = -1$ into the left hand side of the equation we get $0$, whereas we get $17$ on
%the right hand side.  This means $x = -1$ isn't a solution to the algebraic equation.  


\noindent  {\bf Algebraic Example B:}  Is $x = -3$ a solution to equation (\ref{alg_eqn}) above?\\
\vspace{3in}

%However, if we try $x = -3,$ we have
%	\begin{align*}
%	\mbox{LHS}	%&= x^5 + 12x^2 +10x,\\
%								&= (-3)^5 + 12(-3)^2 +10(-3),\\
%								&= -243 + 108 - 30,\\
%								&= -165,\\
%	\end{align*}
%and
%	\begin{align*}
%	\mbox{RHS}	%&= 7x^3 + 24,\\
%									&= 7(-3)^3 + 24,\\
%									&=  -189 + 24,\\
%									&= -165.
%	\end{align*}
%\noindent  Since we get the same result when we substitute $x = -3$ in on the left hand and right hand sides of the
%equation, we can conclude that $x = -3$ is a solution to the algebraic equation.\\
%
\noindent  A couple of notes before we get back to differential equations:
	\begin{itemize}
	\item  At no point did we discuss how we found the solution $x = -3.$  There are multiple ways one might find
		a solution to an algebraic equation (and in a lot of ways figuring out how to solve an equation is far more 
		interesting than knowing a single solution to a single equation),
		but %we don't need to know anything about the solution technique 
		%to be able to verify whether something is or isn't a solution.\\
		\line(1,0){100}\\ \\
		\line(1,0){500}\\ \\
			\line(1,0){500}.\\
		
		One of the first things we will do is ask you to do with a DE 
		to determine whether a given function is or is not a solution.  We are not going to explain (yet)
		how we might have gotten the solution.  We will get there, but an important first step is to understand what it
		means to be a solution to a DE.
		
	\item  It is possible that there are other solutions to the equation.  In fact, there are other solutions to the algebraic
		equation above, and we can see them if we plot the functions 
		$f_1(x) = x^5 + 12x^2 +10x$ and $f_2(x) = 7x^3 + 24$ to see where the curves intersect. 
		
		In other words, %it's possible that an equation has multiple solutions.  
		\line(1,0){400}.\\
		It is also possible that an equation %has no solution (e.g., the equation $x - 1 = x$ has no solution).  
		\line(1,0){300}.\\
		
		We will find that the same is true for DEs; it is possible that a 
		DE has many solutions and it is possible that it has no solution.  For the purposes of this class, this 
		is not terribly important, but we do mention it here so you are aware.\\
	\end{itemize}

\noindent  We have just reviewed what it means to be a solution to an algebraic equation:  {\it if we get the same result 
when we substitute a particular value in on the left hand and right hand sides of an algebraic 
equation, we can conclude that we have a solution to the algebraic equation.}  The definition of solution of a DE
is practically identical:
	
	\begin{quote}  If we get the same result when we substitute a particular {\it function} in on the left hand and right 
	hand sides of a DE, the function is a solution to the DE.  We also describe this situation by saying that the function
	%{\it satisfies the DE.}
	\line(1,0){250}.
	\end{quote}
	
\noindent  Let us look at a couple of examples to see how it works.\\

\noindent  {\bf Example 1:}  Determine if $\ds y = Ce^{x^2} + 3$, where $C$ is any constant, is a solution to the DE 
%$\ds \frac{dy}{dx} = 2xy - 6x.$\\
\line(1,0){100.}

%\noindent  We'll approach this the same way we approached the algebraic equation in the example above:  we will to 
%substitute the function $\ds y = Ce^{x^2} + 3$ in on the left and right hand sides of the equation 
%$\ds \frac{dy}{dx} = 2xy - 6x$ to see
%if we get the same result.
%	\begin{align*}
%	\mbox{LHS}	&= \frac{dy}{dx}\\
%				&= \frac{d}{dx}\Big( Ce^{x^2} + 3 \Big)\\
%				&= C\frac{d}{dx}\Big( e^{x^2} \Big) + \frac{d}{dx}\Big( 3 \Big)\\
%				&= C \cdot e^{x^2} \cdot \frac{d}{dx}\Big( x^2 \Big) + 0\\
%				%margin note:  refer to R/P content on differentiation and chain rule
%				&= C \cdot e^{x^2} \cdot 2x\\
%				&= 2Cxe^{x^2}\\
%	\end{align*}
%	
%	\begin{align*}
%	\mbox{RHS}	&= 2xy - 6x\\
%				&= 2x \Big( Ce^{x^2} + 3 \Big) - 6x\\
%				&= 2Cxe^{x^2} + 6x - 6x\\
%				&= 2Cxe^{x^2}
%	\end{align*}
%We do get the same result when we substitute in on the left and right hand sides of the DE, so we conclude that 
%$\ds y = Ce^{x^2} + 3$ is a solution to the DE $\ds \frac{dy}{dx} = 2xy - 6x.$
\newpage

\noindent  Notice that since we used a generic constant $C$, we have shown that $\ds y = Ce^{x^2} + 3$ is a solution 
for any value we like for the constant $C$.  So, for example, $y = -41e^{x^2} + 3$ is a solution, and so is $y = 199e^{x^2} + 3$.
In other words, the DE $\ds \frac{dy}{dx} = 2xy - 6x$ has an infinite number of solutions.  %mention family of solutions here?
\\

\noindent  {\bf Example 2:}  Determine whether $y = 3\sin(x^2)$ is a solution to the DE %$y' - 12x^2\sin(x^2) = xy''$.
\line(1,0){100}.

%	\begin{align*}
%	y'	&= \frac{d}{dx}\Big( 3\sin(x^2) \Big)\\
%		&= 3\cos(x^2) \cdot \frac{d}{dx}\Big(  x^2\Big) \mbox{\hspace{0.2cm} (using the chain rule)}\\
%		&= 3\cos(x^2) \cdot 2x\\
%		&= 6x\cos(x^2)\\
%		&\\
%	y''	&= 6x \cdot \frac{d}{dx}\Big( \cos(x^2) \Big) + \frac{d}{dx}\Big( 6x \Big)\cdot \cos(x^2)
%				\mbox{\hspace{0.2cm} (using the product rule)}\\
%		&= -12x^2\sin(x^2) + 6\cos(x^2)\\
%	\end{align*}
%	
%\noindent  With these in hand, we can substitute into the left and right hand sides of the equation, as follows.
%	\begin{align*}
%	\mbox{LHS}	&= y' - 12x^2\sin(x^2)\\
%				&= 6x\cos(x^2) - 12x^2\sin(x^2)\\
%				&\\
%	\mbox{RHS}	&= xy''\\
%				&= x\Big( -12x^2\sin(x^2) + 6\cos(x^2) \Big)\\
%				&= -12x^3\sin(x^2) + 6x\cos(x^2)\\
%	\end{align*}
%\noindent  The resultant expressions when we substitute the function into the left hand side and the
%right hand side are similar, but they are not the same.  Hence, $y = 3\sin(x^2)$ is {\bf not a solution}
% to the DE $y' - 12x^2\sin(x^2) = xy''$.\\
\newpage

\noindent {\bf Example 3:}  Determine which of the following are solutions to the DE \line(1,0){100}.
%$y'' - y' - 12y = 0.$
	\begin{align*}
	\mbox{(a)}	\hspace{0.2cm} y	&= 2e^{-3t}\\ %(m+3)(m-4) = m^2 - m - 12
	\mbox{(b)}	\hspace{0.2cm} y	&= e^{4t}\\
	\mbox{(c)}	\hspace{0.2cm} y	&= c_1e^{-3t} + c_2e^{4t}, \mbox{ where $c_1$ and $c_2$ are constants}\\
	\end{align*}
%\noindent  As in the previous example, we need to substitute each function into the DE.  In this case, we need to take
%differentiate each function twice so we are able to substitute in to the DE.
%
%	\begin{align*}
%	\mbox{(a)}	\hspace{0.2cm} y	&= 2e^{-3t}	&\Rightarrow \hspace{0.1cm} 	y'	&= -6e^{-3t}	
%										&\Rightarrow \hspace{0.1cm} y''	&= 18e^{-3t}\\
%							&			&							&	&&\\
%	\mbox{(b)}	\hspace{0.2cm} y	&= e^{4t}		&\Rightarrow \hspace{0.1cm}	 y'	&= 4e^{4t}		
%										&\Rightarrow \hspace{0.1cm} y''	&= 16e^{4t}\\
%							&			&							&	&&\\
%	\mbox{(c)}	\hspace{0.2cm} y	&= c_1e^{-3t} + c_2e^{4t} 
%										&\Rightarrow \hspace{0.1cm} 	y'	&= -3c_1e^{-3t} + 4c_2e^{4t}
%										&\Rightarrow \hspace{0.1cm} y''	&= 9c_1e^{-3t} + 16c_2e^{4t}
%	\end{align*}
%\noindent  Now let's check each provided function by substituting into the DE.
%	\begin{align*}
%	\mbox{(a)}	\hspace{0.2cm} LHS	&= y'' - y' - 12y\\
%							&= \Big(18e^{-3t}\Big) - \Big( -6e^{-3t} \Big) - 12\Big( 2e^{-3t} \Big)\\
%							&= 18e^{-3t} + 6e^{-3t} - 24e^{-3t}\\
%							&= 0\\
%							&= RHS\\
%	\end{align*}
%
%	\begin{align*}
%	\mbox{(a)}	\hspace{0.2cm} LHS	&= y'' - y' - 12y\\
%							&= \Big(16e^{4t}\Big) - \Big( 4e^{4t} \Big) - 12\Big( e^{4t} \Big)\\
%							&= 16e^{4t} - 4e^{4t} - 12e^{4t}\\
%							&= 0\\
%							&= RHS\\
%	\end{align*}
%	
%	\begin{align*}
%	\mbox{(c)}	\hspace{0.2cm} LHS	&= y'' - y' - 12y\\
%							&= \Big(9c_1e^{-3t} + 16c_2e^{4t}\Big) - \Big( -3c_1e^{-3t} + 4c_2e^{4t} \Big) - 12\Big( c_1e^{-3t} + c_2e^{4t} \Big)\\
%							&= 9c_1e^{-3t} + 16c_2e^{4t} + 3c_1e^{-3t} - 4c_2e^{4t} - 12c_1e^{-3t} - 12c_2e^{4t}\\
%							&= \Big( 9c_1e^{-3t} + 3c_1e^{-3t} - 12c_1e^{-3t}   \Big) + \Big( 16c_2e^{4t} - 4c_2e^{4t} - 12c_2e^{4t} \Big)\\
%							& = 0\\
%							&= RHS\\
%	\end{align*}
%
%\noindent  It turns out that each of the functions provided is, in fact, a solution to the DE.\\


\end{document}